% BHU PYQ -- Professional Exam Template
\documentclass[11pt,a4paper]{article}
\usepackage[a4paper,top=2.5cm,bottom=2.5cm,left=2.8cm,right=2.8cm,headheight=14pt]{geometry}
\usepackage{amsmath,amssymb}
\usepackage[T1]{fontenc}
\usepackage[utf8]{inputenc}
\usepackage{lmodern}
\usepackage{microtype}
\usepackage{fancyhdr}
\usepackage{enumitem}
\usepackage{hyperref}
\usepackage{lastpage}
\usepackage{parskip}

\hypersetup{colorlinks=true,urlcolor=black,linkcolor=black,
  pdftitle={CSCMJ21/CSCMN21: Data Structures and Algorithms I -- BHU PYQ 2024-25}}

\pagestyle{fancy}\fancyhf{}
\renewcommand{\headrulewidth}{0.4pt}\renewcommand{\footrulewidth}{0.4pt}
\fancyhead[L]{\small\textit{Banaras Hindu University}}
\fancyhead[R]{\small\textit{CSCMJ21/CSCMN21: Data Structures \& Algorithms I}}
\fancyfoot[C]{\small Page \thepage\ of \pageref{LastPage}}

\newlist{parts}{enumerate}{2}
\setlist[parts,1]{label=(\alph*),leftmargin=*,itemsep=4pt,topsep=3pt}
\setlist[parts,2]{label=(\roman*),leftmargin=*,itemsep=2pt,topsep=2pt}
\newcommand{\pts}[1]{\hfill\textit{[#1]}}

\begin{document}
\begin{center}
  {\large\bfseries BANARAS HINDU UNIVERSITY}\[2pt]
  {\normalsize B.Sc.\ (Hons.) Semester II Examination (NEP), 2024-25}\[6pt]
  \rule{0.8\linewidth}{0.5pt}\[6pt]
  {\large\bfseries COMPUTER SCIENCE}\[4pt]
  {\large\bfseries Paper: CSCMJ21 / CSCMN21 --- Data Structures and Algorithms I}\[4pt]
  {\normalsize Time: 2:30 Hours \hfill Full Marks: 70}
\end{center}
\rule{\linewidth}{0.4pt}
\smallskip
\noindent\textit{Instructions: (i) The question paper contains 7 questions out of which you are required to answer any 5 questions. The question paper is of 70 marks with each question carrying 14 marks.\
(ii) Answer to each question should start from a fresh page. Write your Roll No.\ at the top immediately on receipt of this question paper.}
\bigskip

\noindent\textbf{Question 1.}
\begin{parts}
  \item Explain linear and non-linear data structures with suitable examples.\pts{7}
  \item Define queue, priority queue, circular queue and double-ended queue with suitable illustrations.\pts{7}
\end{parts}

\bigskip
\noindent\textbf{Question 2.}
\begin{parts}
  \item Explain hashing and hash table with suitable example. Discuss collision resolution techniques with illustrations.\pts{10}
  \item Explain load factor in hashing with suitable example.\pts{4}
\end{parts}

\bigskip
\noindent\textbf{Question 3.}
\begin{parts}
  \item Explain insertion sort with suitable example. Sort the following elements using quick sort:\pts{14}
  \[85, 09, 101, 01, 08, 99, 80, 76\]
\end{parts}

\bigskip
\noindent\textbf{Question 4.}
\begin{parts}
  \item Explain Divide and Conquer and Dynamic Programming approach with suitable example. Discuss Graph as a data structure.\pts{14}
\end{parts}

\bigskip
\noindent\textbf{Question 5.}
\begin{parts}
  \item Define asymptotic notations through an example.\pts{7}
  \item Explain balanced binary tree and explain how to check it.\pts{7}
\end{parts}

\bigskip
\noindent\textbf{Question 6.}
\begin{parts}
  \item Define shortest path problem with illustrations.\pts{7}
  \item Explain spanning tree and minimum spanning tree. Give suitable illustrations.\pts{7}
\end{parts}

\bigskip
\noindent\textbf{Question 7.}
\noindent Write brief notes on any \textbf{two} of the following:\pts{14}
\begin{parts}
  \item Knapsack Problem
  \item Tree Traversal Algorithms
  \item Binary Search Tree
  \item Adjacency List and Adjacency Matrix
\end{parts}

\vfill\rule{\linewidth}{0.4pt}
\begin{center}\small
  \href{https://www.sciqb.com/}{Click here to visit our website}\[4pt]
  \href{https://me-aryan.vercel.app/}{Click here to visit our contributor}\[4pt]
  \href{https://quashan-me.vercel.app/}{Click here to visit our contributor}
\end{center}
\end{document}